\documentclass[a4paper,11pt]{article}

% --- Geometry, layout ---
\usepackage[a4paper, margin=2.5cm]{geometry}
\usepackage{parskip}
\usepackage[none]{hyphenat}

% --- Maths ---
\usepackage{amsmath}
\usepackage{amssymb}
\usepackage{amsthm}
\usepackage{mathtools}
\usepackage{bm}

% --- Floats and tables ---
\usepackage{graphicx}
\usepackage{float}
\usepackage{booktabs}
\usepackage{siunitx}
\sisetup{
  table-number-alignment = center,
  exponent-product = \cdot,
  output-decimal-marker = {.},
  table-format = 1.4e-1,
}
\usepackage{subcaption}

% --- Code listings ---
\usepackage{listings}
\usepackage{xcolor}
\lstset{
  basicstyle      = \ttfamily\small,
  keywordstyle    = \color{blue!70!black}\bfseries,
  commentstyle    = \color{gray}\itshape,
  stringstyle     = \color{green!50!black},
  showstringspaces= false,
  breaklines      = true,
  frame           = single,
  rulecolor       = \color{gray!50},
  framesep        = 4pt,
}

% --- Bibliography ---
\usepackage{natbib}

% --- Cross-references and hyperlinks (load hyperref before cleveref) ---
\usepackage[colorlinks=true,
            linkcolor=blue!60!black,
            citecolor=blue!60!black,
            urlcolor=blue!60!black]{hyperref}
\usepackage{cleveref}

% --- Operators and macros consistent with Schwedes et al. (2017) ---
\newcommand{\R}{\mathbb{R}}
\newcommand{\Om}{\Omega}
\newcommand{\dom}{\partial\Omega}
\newcommand{\Lp}[1]{L^{2}(#1)}              % L^2 space
\newcommand{\inner}[3]{(#1,\,#2)_{#3}}        % inner product
\newcommand{\dx}{\,\mathrm{d}x}
\newcommand{\ds}{\,\mathrm{d}s}
\newcommand{\dt}{\,\mathrm{d}t}
\newcommand{\Rmap}{\mathcal{R}}              % Riesz map (Schwedes notation)
\newcommand{\Jt}{\tilde{J}}                  % reduced functional
\newcommand{\Vstar}{V^{*}}                   % dual space
\newcommand{\Mspace}{M}
\newcommand{\Vspace}{V}
\newcommand{\Vtest}{\mathcal{V}}
\DeclareMathOperator{\grad}{\nabla}

% --- Theorem environments ---
\newtheorem{counter}{Counter}[section]

\newtheoremstyle{spacedplain}
  {\parskip}{\parskip}{\itshape}{}{\bfseries}{.}{.5em}{}
\newtheoremstyle{spaceddefinition}
  {\parskip}{\parskip}{}{}{\bfseries}{.}{.5em}{}

\theoremstyle{spacedplain}
\newtheorem{theorem}[counter]{Theorem}
\newtheorem{lemma}[counter]{Lemma}
\newtheorem{prop}[counter]{Proposition}
\newtheorem{cor}[counter]{Corollary}

\theoremstyle{spaceddefinition}
\newtheorem{defn}[counter]{Definition}
\newtheorem{example}[counter]{Example}
\newtheorem{remark}[counter]{Remark}

% --- Cleveref names ---
\crefname{section}{Section}{Sections}
\crefname{subsection}{Section}{Sections}
\crefname{equation}{equation}{equations}
\crefname{table}{Table}{Tables}
\crefname{figure}{Figure}{Figures}
\crefname{defn}{Definition}{Definitions}
\crefname{theorem}{Theorem}{Theorems}
\crefname{prop}{Proposition}{Propositions}

% --- Title block ---
\title{Optimal Control of the Poisson and Burgers Equations,\\
       with a Study of Mesh-Dependence in the Optimisation}
\author{Brian Chang}
\date{\today}

% =============================================================================
\begin{document}
% =============================================================================

\maketitle
\tableofcontents
\newpage

% =============================================================================
\section{Introduction}
\label{sec:intro}
% =============================================================================

Many problems arising in science and engineering involve determining
unknown inputs to a physical system from partial or indirect observations
of its output. Such problems are commonly formulated as PDE-constrained
optimisation problems, in which the behaviour of the system is governed
by a partial differential equation, while the unknown input is determined
by minimising a suitable objective functional.

In this work, we consider two model problems that together exemplify
the linear and the nonlinear cases of PDE-constrained optimisation. The
first is the steady Poisson equation on a bounded domain $\Om \subset \R^d$,
in which we seek a source term acting within the domain that drives the
solution towards a prescribed desired state. The second is the
time-dependent scalar viscous Burgers equation, which adds the new
features of nonlinearity and a time-dependent (and backward-in-time)
adjoint solve.

After deriving the weak forms, the adjoint equations and the
$L^2$-gradients of the reduced functionals for these two problems, we
examine the impact that the choice of inner product on the control space
has on the numerical optimisation. As discussed by
\citet[Sect.~2.1]{shfp2017}, many widely used optimisation packages with
Python interfaces (including \texttt{scipy.optimize} and \texttt{TAO})
hard-code the Euclidean $\ell^2$ inner product on the coefficient
vector. The Riesz representation theorem identifies derivatives in the
dual space with primal elements only relative to the inner product
placed on the control space; the resulting gradient is
\emph{mesh-dependent} when the wrong inner product is used, in the sense
that the iteration count of a quasi-Newton method grows with mesh
refinement. We reproduce, and extend beyond, Table~2.2 of
\citet{shfp2017} to demonstrate this effect on the Poisson optimal
control problem.

The remainder of this work is organised as follows.
\Cref{sec:preliminaries} establishes the functional-analytic
framework, including the calculus of Fr{\'e}chet derivatives and the
Riesz representation theorem, which we will need throughout.
\Cref{sec:poisson-weak} derives the weak form of the Poisson forward
problem and \cref{sec:poisson-optimisation} sets up the corresponding
optimisation problem. The derivation of the tangent linear and adjoint
equations, and of the $L^2$-gradient of the reduced functional, occupies
\cref{sec:poisson-derivative}.
\Cref{sec:poisson-numerics} verifies the resulting numerical pipeline
on the linear problem. We then treat the scalar viscous Burgers problem
in \cref{sec:burgers-weak,sec:burgers-derivative,sec:burgers-numerics}.
\Cref{sec:mesh-dependence} contains the mesh-dependence study.

\vspace{0.5em}
\noindent
Throughout the work, references in the form
\textit{``\citep[p.~XX]{shfp2017}''} or
\textit{``\citep[Sect.~Y.Z]{shfp2017}''} point to the reference book
\citep{shfp2017}, which we follow closely both in notation and in
structure.


% =============================================================================
\section{Mathematical Preliminaries}
\label{sec:preliminaries}
% =============================================================================

We introduce the functional-analytic objects used throughout the work.
The material follows the standard treatment in Chapter~1 of
\citet{shfp2017}; we collect it here so that the calculus of variations
in function spaces, including the Fr{\'e}chet derivative and the Riesz
representation theorem, is in hand by the time we derive the adjoint
equation.

\subsection{Domain and notation}
\label{sec:prelim:domain}

Let $\Om \subset \R^d$, with $d \in \{1, 2, 3\}$, be a bounded open
domain with Lipschitz boundary $\dom$. Unless stated otherwise, all
functions are defined on $\Om$.

We denote by $\Lp{\Om}$ the space of square-integrable functions on
$\Om$, equipped with the inner product and norm
\begin{equation}
  \inner{u}{v}{\Lp{\Om}} := \int_{\Om} u v \dx,
  \qquad
  \|u\|_{\Lp{\Om}} := \inner{u}{u}{\Lp{\Om}}^{1/2}.
\end{equation}

\subsection{Sobolev spaces}
\label{sec:prelim:sobolev}

\begin{defn}[Sobolev space \texorpdfstring{$H^1(\Om)$}{H1}]
The Sobolev space $H^1(\Om)$ is defined by
\begin{equation}
  H^1(\Om) := \left\{ u \in \Lp{\Om} \;\middle|\;
                      \grad u \in (\Lp{\Om})^d \right\}.
\end{equation}
\end{defn}

A function $u$ belongs to $H^1(\Om)$ if both the function itself and
its first (weak) derivatives are square-integrable. The space is
equipped with the norm
\begin{equation}
  \|u\|_{H^1(\Om)}^2
    := \|u\|_{\Lp{\Om}}^2 + \|\grad u\|_{\Lp{\Om}}^2.
\end{equation}
To incorporate homogeneous Dirichlet boundary conditions into the
function space itself, we introduce the subspace $H_0^1(\Om)$.

\begin{defn}[Sobolev space \texorpdfstring{$H_0^1(\Om)$}{H10}]
The space $H_0^1(\Om)$ is defined as the closure of $C_c^\infty(\Om)$
in the $H^1(\Om)$ norm.
\end{defn}

Here, $C_c^\infty(\Om)$ denotes the smooth functions with compact
support in $\Om$, that is, smooth functions that vanish in a
neighbourhood of $\dom$. Functions in $H_0^1(\Om)$ have zero boundary
values in the sense of Sobolev traces; $H_0^1(\Om)$ is therefore the
natural choice of solution space for problems with homogeneous
Dirichlet boundary conditions.

\subsection{Dual spaces and bilinear forms}
\label{sec:prelim:duals}

The dual space of $H_0^1(\Om)$ is denoted by $H^{-1}(\Om)$ and
consists of bounded linear functionals on $H_0^1(\Om)$. Duality
pairings between $H^{-1}(\Om)$ and $H_0^1(\Om)$ are written as
$\langle \cdot, \cdot \rangle$. When no confusion can arise, inner
products in Hilbert spaces are written using parentheses, and duality
pairings using angle brackets.

Let $V$ be a Hilbert space. A bilinear form
$a(\cdot, \cdot) : V \times V \to \R$ is \emph{continuous} if there
exists a constant $C > 0$ such that
\begin{equation}
  |a(u, v)| \le C \|u\|_V \|v\|_V
  \qquad \forall u, v \in V,
\end{equation}
and \emph{coercive} if there exists a constant $\alpha > 0$ such that
\begin{equation}
  a(u, u) \ge \alpha \|u\|_V^2
  \qquad \forall u \in V.
\end{equation}
These two properties are sufficient hypotheses of the Lax-Milgram
theorem, which guarantees existence and uniqueness of weak solutions
for a wide class of linear elliptic problems.

\subsection{Control space}
\label{sec:prelim:control}

In addition to the state variable, the problems considered in this
work involve a control variable. The control belongs to a Hilbert
space $\Mspace$, equipped with inner product
$\inner{\cdot}{\cdot}{\Mspace}$ and associated norm $\|\cdot\|_{\Mspace}$.

For the Poisson model problem of \cref{sec:poisson-optimisation} we
take $\Mspace = \Lp{\Om}$. This choice lets us treat the source term in
the state equation naturally within the variational framework.

\subsection{Differential calculus in function spaces}
\label{sec:prelim:frechet}

Formulating PDE-constrained optimisation problems requires a notion of
differentiation for operators and functionals mapping between Hilbert
spaces. Let $U$ and $V$ be Hilbert spaces, and let
$F : U \to V$ be a (possibly nonlinear) operator.

\begin{defn}[Fr{\'e}chet derivative]
\label{defn:frechet}
The operator $F$ is \emph{Fr{\'e}chet differentiable} at $u \in U$
if there exists a bounded linear operator
$\mathrm{d} F(u) \in \mathcal{L}(U, V)$ such that
\begin{equation}
  \lim_{\|\delta u\|_U \to 0}
    \frac{\|F(u + \delta u) - F(u) - \mathrm{d} F(u)\, \delta u\|_V}
         {\|\delta u\|_U}
  = 0.
\end{equation}
The directional derivative in the direction $\delta u \in U$ is
written $\mathrm{d} F(u; \delta u) := \mathrm{d} F(u)\, \delta u$.
\end{defn}

When an operator or functional depends on multiple variables, the
notion of Fr{\'e}chet differentiation extends naturally to partial
derivatives.

\begin{defn}[Partial Fr{\'e}chet derivative]
\label{defn:partial-frechet}
Let $F : U \times \Mspace \to V$. The \emph{partial Fr{\'e}chet
derivative} of $F$ with respect to $u$ at the point $(u, m)$, acting
in the direction $\delta u \in U$, is denoted
$\partial_u F(u, m; \delta u) \in V$. It is defined as the Fr{\'e}chet
derivative of the map $u \mapsto F(u, m)$ with $m$ held fixed. The
partial derivative with respect to $m$, denoted
$\partial_m F(u, m; \delta m)$, is defined analogously.
\end{defn}

Finally, the derivative of a composite functional is obtained from the
chain rule.

\begin{theorem}[Chain rule]
\label{thm:chain}
Let $\Mspace$ and $U$ be Hilbert spaces. Let $S : \Mspace \to U$ be
Fr{\'e}chet differentiable at $m \in \Mspace$, and let
$J : U \times \Mspace \to \R$ be continuously Fr{\'e}chet
differentiable at $(S(m), m)$. Then the reduced functional
$\Jt(m) := J(S(m), m)$ is Fr{\'e}chet differentiable at $m$, and its
derivative in the direction $\delta m \in \Mspace$ is given by
\begin{equation}
  \mathrm{d} \Jt(m;\, \delta m)
  = \partial_u J\bigl( S(m), m;\, \mathrm{d} S(m;\, \delta m) \bigr)
  + \partial_m J\bigl( S(m), m;\, \delta m \bigr).
\end{equation}
\end{theorem}

\subsection{The Riesz representation theorem}
\label{sec:prelim:riesz}

To properly define a gradient of a functional on a Hilbert space, we
must understand the relationship between the space and its dual. Let
$H$ be a Hilbert space with inner product
$\inner{\cdot}{\cdot}{H}$.

\begin{defn}[Dual space]
The dual space of $H$, denoted $H^{*}$, is the space of all continuous
linear functionals on $H$ \citep[p.~21]{shfp2017}:
\begin{equation}
  H^{*} := \bigl\{ f : H \to \R \;\big|\;
                    f \text{ is linear and continuous} \bigr\}.
\end{equation}
\end{defn}

The Hilbert space $H$ and its dual $H^{*}$ are isometrically
isomorphic. The Riesz representation theorem makes this precise.

\begin{theorem}[Riesz representation theorem]
\label{thm:riesz}
Let $H$ be a Hilbert space with inner product
$\inner{\cdot}{\cdot}{H}$. For any continuous linear functional
$f \in H^{*}$, there exists a unique $u_f \in H$ such that
\begin{equation}
  f(v) = \inner{u_f}{v}{H}
  \qquad \forall v \in H,
\end{equation}
and the norms agree, $\|f\|_{H^{*}} = \|u_f\|_H$.
\end{theorem}

The theorem identifies each continuous linear functional with a unique
element of $H$. This identification is the \emph{Riesz map}.

\begin{defn}[Riesz map and Riesz representer]
\label{defn:rieszmap}
The map $\Rmap_H : H^{*} \to H$ that sends each $f \in H^{*}$ to its
unique representative $u_f \in H$ is called the \emph{Riesz map}.
The element $\Rmap_H(f) := u_f$ is the \emph{Riesz representer} of $f$
in $H$.
\end{defn}

The Riesz map depends on the choice of inner product on $H$: the same
linear functional has different Riesz representers under different
inner products. This dependence becomes a central concern in
\cref{sec:mesh-dependence}, where we show that using the wrong Riesz
map yields a mesh-dependent gradient and hence a mesh-dependent
optimisation algorithm.

\subsection{Weak formulations}
\label{sec:prelim:weakforms}

A weak formulation of a partial differential equation is obtained by
multiplying the strong form by a test function, integrating over the
domain, and applying integration by parts. The transfer of derivatives
from the solution onto the test functions lets us seek solutions in
Sobolev spaces rather than in spaces of smooth functions. The same
machinery underlies the finite element approximation studied in
\cref{sec:poisson-numerics}.


% =============================================================================
\section{The Forward Poisson Problem}
\label{sec:poisson-weak}
% =============================================================================

We now derive the weak formulation of the Poisson equation that defines
the state variable in the model problem. Our goal is to rewrite the PDE
in a form that involves only first derivatives of the solution, so that
we can seek solutions in $H_0^1(\Om)$ rather than in spaces of
classical (twice differentiable) solutions.

\subsection{Strong formulation}
\label{sec:poisson-weak:strong}

Let $m$ be a given source term (for example $m \in \Lp{\Om}$). The
strong form of the forward problem is
\begin{equation}
  \label{eq:strong-forward}
  -\Delta u = m \quad \text{in } \Om,
  \qquad
  u = 0 \quad \text{on } \dom.
\end{equation}

\subsection{Derivation of the weak form}
\label{sec:poisson-weak:derivation}

We choose a test function in
\begin{equation}
  v \in H_0^1(\Om),
\end{equation}
which vanishes on $\dom$ in the trace sense. Multiplying
\cref{eq:strong-forward} by $v$ and integrating over $\Om$ yields
\begin{equation}
  \label{eq:weak-derivation-step1}
  \int_{\Om} (-\Delta u)\, v \dx = \int_{\Om} m\, v \dx.
\end{equation}
The left-hand side contains second derivatives of $u$. We remove one
of them by applying Green's first identity, which gives
\begin{equation}
  \int_{\Om} (-\Delta u)\, v \dx
  = \int_{\Om} \grad u \cdot \grad v \dx
  - \int_{\dom} \frac{\partial u}{\partial n}\, v \ds,
\end{equation}
where $n$ denotes the outward unit normal on $\dom$ and
$\partial u / \partial n := \grad u \cdot n$ is the normal derivative of
$u$. Because $v \in H_0^1(\Om)$ has zero trace on $\dom$, the boundary
integral vanishes, and we obtain
\begin{equation}
  \int_{\Om} (-\Delta u)\, v \dx
  = \int_{\Om} \grad u \cdot \grad v \dx.
\end{equation}
Substituting this into \cref{eq:weak-derivation-step1} gives
\begin{equation}
  \int_{\Om} \grad u \cdot \grad v \dx
  = \int_{\Om} m\, v \dx,
\end{equation}
and this identity must hold for all test functions $v \in H_0^1(\Om)$.
We collect the result into the following definition.

\begin{defn}[Weak formulation of the forward Poisson problem]
\label{defn:weak-poisson}
Let $m \in \Lp{\Om}$. The weak formulation of \cref{eq:strong-forward}
is: find $u \in H_0^1(\Om)$ such that
\begin{equation}
  \label{eq:weak-forward}
  \int_{\Om} \grad u \cdot \grad v \dx
  = \int_{\Om} m\, v \dx
  \qquad \forall v \in H_0^1(\Om).
\end{equation}
\end{defn}

It is convenient to introduce the bilinear form
\begin{equation}
  \label{eq:bilinear-b}
  b(u, v) := \int_{\Om} \grad u \cdot \grad v \dx
\end{equation}
and the linear functional
\begin{equation}
  \label{eq:linear-ell}
  \ell(v) := \int_{\Om} m\, v \dx,
\end{equation}
so that \cref{eq:weak-forward} reads
\begin{equation}
  \label{eq:weak-forward-compact}
  \text{find } u \in H_0^1(\Om) \text{ such that }
  b(u, v) = \ell(v)
  \quad \forall v \in H_0^1(\Om).
\end{equation}
The bilinear form $b$ is continuous and, by the Poincar{\'e}
inequality, coercive on $H_0^1(\Om)$. The linear functional $\ell$ is
bounded on $H_0^1(\Om)$ whenever $m \in \Lp{\Om}$. The Lax-Milgram
theorem then guarantees existence and uniqueness of the weak solution.


% =============================================================================
\section{The Poisson Optimal Control Problem}
\label{sec:poisson-optimisation}
% =============================================================================

We now introduce the PDE-constrained optimisation problem studied in
the linear setting. We seek a source term $m$ such that the
corresponding solution $u$ of \cref{eq:strong-forward} matches a given
desired state $d$, while penalising large controls.

Throughout this and the following sections, we set
\begin{equation}
  \label{eq:VM-spaces}
  V := H_0^1(\Om),
  \qquad
  \Mspace := \Lp{\Om}.
\end{equation}

\subsection{Objective functional}
\label{sec:opt:objective}

Let $d \in \Lp{\Om}$ be a given desired state and let $\alpha > 0$.

\begin{defn}[Objective functional]
\label{defn:objective}
The objective functional $J : V \times \Mspace \to \R$ is defined by
\begin{equation}
  \label{eq:objective-J}
  J(u, m)
  := \tfrac{1}{2} \|u - d\|_{\Lp{\Om}}^2
  + \tfrac{\alpha}{2} \|m\|_{\Mspace}^2.
\end{equation}
\end{defn}

The first term measures the misfit between the state $u$ and the
desired state $d$. The second term penalises large controls; the
parameter $\alpha > 0$ balances accuracy of fit against regularity of
the control. The presence of the regularisation term ensures that the
optimisation problem is well posed.

\subsection{The constrained optimisation problem}
\label{sec:opt:constrained}

Combining \cref{defn:objective} with the weak state equation
\cref{eq:weak-forward-compact}, the PDE-constrained optimisation
problem reads as follows.

\begin{defn}[PDE-constrained optimisation problem]
\label{defn:full-problem}
Find $(u, m) \in V \times \Mspace$ minimising
\begin{equation}
  \min_{(u, m) \in V \times \Mspace} J(u, m)
\end{equation}
subject to
\begin{equation}
  b(u, v) = \ell(v) \qquad \forall v \in V,
\end{equation}
where the bilinear form $b$ and the linear functional $\ell$ are
those of \cref{eq:bilinear-b,eq:linear-ell}.
\end{defn}

\subsection{Reduced formulation}
\label{sec:opt:reduced}

For each $m \in \Mspace$ there exists a unique state $u \in V$
satisfying the state equation, so we may view the state as a function
of the control.

\begin{defn}[Solution operator]
\label{defn:solution-operator}
The solution operator $S : \Mspace \to V$ assigns to each
$m \in \Mspace$ the unique solution $S(m) := u(m) \in V$ of the weak
state equation associated with $m$.
\end{defn}

Using the solution operator we eliminate the state from the
optimisation problem.

\begin{defn}[Reduced functional]
\label{defn:reduced-functional}
The reduced functional $\Jt : \Mspace \to \R$ is defined by
\begin{equation}
  \Jt(m) := J(S(m), m) = J(u(m), m).
\end{equation}
\end{defn}

The constrained problem of \cref{defn:full-problem} is then equivalent
to the unconstrained minimisation problem
\begin{equation}
  \min_{m \in \Mspace} \Jt(m).
\end{equation}


% =============================================================================
\section{Derivative of the Reduced Functional and the Adjoint Equation}
\label{sec:poisson-derivative}
% =============================================================================

We derive an expression for the Fr{\'e}chet derivative of the reduced
functional $\Jt$ via the adjoint method. The derivation uses the chain
rule (\cref{thm:chain}) and the calculus of partial Fr{\'e}chet
derivatives (\cref{defn:partial-frechet}) introduced in
\cref{sec:prelim:frechet}.

Applying the chain rule to $\Jt(m) = J(u(m), m)$ in the direction
$\delta m \in \Mspace$ gives
\begin{equation}
  \label{eq:chain-rule-applied}
  \mathrm{d}\Jt(m;\, \delta m)
  = \partial_u J(u, m;\, \mathrm{d} u_m(m;\, \delta m))
  + \partial_m J(u, m;\, \delta m),
\end{equation}
where $\mathrm{d} u_m(m;\, \delta m)$ is the total derivative of the
state with respect to the control in the direction $\delta m$
\citep[p.~32]{shfp2017}. Writing
\begin{equation}
  \delta u := \mathrm{d} u_m(m;\, \delta m) \in V,
\end{equation}
this becomes
\begin{equation}
  \label{eq:chain-rule-shorter}
  \mathrm{d}\Jt(m;\, \delta m)
  = \partial_u J(u, m;\, \delta u)
  + \partial_m J(u, m;\, \delta m).
\end{equation}

\subsection{Tangent linear equation}
\label{sec:deriv:tangent}

To compute $\delta u$ we differentiate the state equation. We define a
\emph{forward model operator} that encodes the residual of the weak
problem.

\begin{defn}[Forward model operator]
\label{defn:forward-operator}
The forward model operator is defined as the residual of the weak
state equation:
\begin{equation}
  F(u, m;\, v) := b(u, v) - \inner{m}{v}{\Lp{\Om}}.
\end{equation}
\end{defn}

The weak state equation can then be written as
$F(u(m), m;\, v) = 0$ for all $v \in V$ \citep[p.~29]{shfp2017}.
Differentiating with respect to $m$, the chain rule yields
\begin{equation}
  \label{eq:tlm-implicit}
  \partial_u F(u, m;\, \delta u, v)
  + \partial_m F(u, m;\, \delta m, v) = 0
  \qquad \forall v \in V.
\end{equation}
We compute the two partial derivatives.

\paragraph{Derivative with respect to $u$.}
For $\varepsilon \in \R$, set $u_\varepsilon := u + \varepsilon \delta u$.
Using linearity of $b$ in its first argument,
\begin{equation}
  F(u_\varepsilon, m;\, v)
  = b(u, v) - \inner{m}{v}{\Lp{\Om}}
  + \varepsilon\, b(\delta u, v).
\end{equation}
The corresponding difference quotient is
\begin{equation}
  \frac{F(u + \varepsilon \delta u, m;\, v) - F(u, m;\, v)}{\varepsilon}
  = b(\delta u, v).
\end{equation}
Taking $\varepsilon \to 0$ gives
\begin{equation}
  \label{eq:partial-Fu}
  \partial_u F(u, m;\, \delta u, v) = b(\delta u, v).
\end{equation}

\paragraph{Derivative with respect to $m$.}
By the analogous argument with $m_\varepsilon := m + \varepsilon \delta m$,
\begin{equation}
  \label{eq:partial-Fm}
  \partial_m F(u, m;\, \delta m, v)
  = -\inner{\delta m}{v}{\Lp{\Om}}.
\end{equation}

Combining \cref{eq:partial-Fu,eq:partial-Fm} with
\cref{eq:tlm-implicit} gives the \emph{tangent linear equation}, also
known as the tangent linear model \citep[p.~32]{shfp2017}.

\begin{prop}[Tangent linear equation]
\label{prop:tangent-linear}
The state variation $\delta u \in V$ satisfies
\begin{equation}
  \label{eq:tangent-linear}
  b(\delta u, v) = \inner{\delta m}{v}{\Lp{\Om}}
  \qquad \forall v \in V.
\end{equation}
\end{prop}

\subsection{Partial derivatives of the objective functional}
\label{sec:deriv:Jpartial}

\paragraph{Derivative of $J$ with respect to $u$.}
Setting $u_\varepsilon := u + \varepsilon v$ and expanding,
\begin{equation}
  J(u_\varepsilon, m)
  = J(u, m)
  + \varepsilon \inner{u - d}{v}{\Lp{\Om}}
  + \tfrac{\varepsilon^2}{2} \|v\|_{\Lp{\Om}}^2.
\end{equation}
Therefore
\begin{equation}
  \label{eq:partial-Ju}
  \partial_u J(u, m;\, v) = \inner{u - d}{v}{\Lp{\Om}}.
\end{equation}

\paragraph{Derivative of $J$ with respect to $m$.}
Setting $m_\varepsilon := m + \varepsilon \delta m$ and expanding
$\|m + \varepsilon \delta m\|^2$,
\begin{equation}
  J(u, m_\varepsilon)
  = J(u, m)
  + \alpha \varepsilon \inner{m}{\delta m}{\Lp{\Om}}
  + \tfrac{\alpha \varepsilon^2}{2} \|\delta m\|_{\Lp{\Om}}^2.
\end{equation}
Therefore
\begin{equation}
  \label{eq:partial-Jm}
  \partial_m J(u, m;\, \delta m) = \alpha \inner{m}{\delta m}{\Lp{\Om}}.
\end{equation}

\subsection{Adjoint problem}
\label{sec:deriv:adjoint}

To evaluate the derivative of $\Jt$ without solving the tangent linear
equation for every direction $\delta m$, we introduce an adjoint
variable \citep[p.~33]{shfp2017}.

\begin{defn}[Adjoint variable]
\label{defn:adjoint}
The adjoint variable $\lambda \in V$ is the unique solution of
\begin{equation}
  \partial_u F(u, m;\, \lambda, v)
  = \partial_u J(u, m;\, v)
  \qquad \forall v \in V.
\end{equation}
\end{defn}

Substituting \cref{eq:partial-Fu,eq:partial-Ju} into the definition
yields
\begin{equation}
  \label{eq:adjoint-implicit}
  b(\lambda, v) = \inner{u - d}{v}{\Lp{\Om}}
  \qquad \forall v \in V.
\end{equation}
Because $b$ is symmetric for the Poisson problem, this is equivalent to
\begin{equation}
  \label{eq:adjoint-symmetric}
  b(v, \lambda) = \inner{u - d}{v}{\Lp{\Om}}
  \qquad \forall v \in V.
\end{equation}
The bilinear form $b$ is continuous and coercive on $V$, and the
right-hand side is a continuous linear functional on $V$, so
existence and uniqueness of the adjoint follow from Lax-Milgram
\citep[p.~14]{shfp2017}.

\subsection{Derivative of the reduced functional}
\label{sec:deriv:reduced}

Substituting \cref{eq:partial-Ju,eq:partial-Jm} into the chain-rule
expression \cref{eq:chain-rule-shorter} gives
\begin{equation}
  \mathrm{d}\Jt(m;\, \delta m)
  = \inner{u - d}{\delta u}{\Lp{\Om}}
  + \alpha \inner{m}{\delta m}{\Lp{\Om}}.
\end{equation}
We now eliminate $\delta u$ using the adjoint and tangent linear
equations. Choosing $v = \delta u$ in
\cref{eq:adjoint-symmetric},
\begin{equation}
  \inner{u - d}{\delta u}{\Lp{\Om}} = b(\delta u, \lambda).
\end{equation}
Choosing $v = \lambda$ in \cref{eq:tangent-linear},
\begin{equation}
  b(\delta u, \lambda) = \inner{\delta m}{\lambda}{\Lp{\Om}}.
\end{equation}
Hence
\begin{equation}
  \label{eq:dJtilde-pre-riesz}
  \mathrm{d}\Jt(m;\, \delta m)
  = \inner{\lambda}{\delta m}{\Lp{\Om}}
  + \alpha \inner{m}{\delta m}{\Lp{\Om}}
  \qquad \forall \delta m \in \Mspace,
\end{equation}
where $u = u(m)$ is the state corresponding to $m$ and $\lambda$ is
the adjoint solution.

\subsection{The \texorpdfstring{$L^2$}{L2}-gradient}
\label{sec:deriv:l2grad}

The reduced derivative \cref{eq:dJtilde-pre-riesz} is a continuous
linear functional acting on directions $\delta m \in \Mspace$, hence
an element of the dual space $\Mspace^{*}$. By the Riesz
representation theorem (\cref{thm:riesz}), there exists a unique
element of $\Mspace$ representing it.

\begin{defn}[$L^2$-gradient]
\label{defn:l2gradient}
Let $\Rmap_{L^2} : \Mspace^{*} \to \Mspace$ denote the Riesz map
associated with the $\Lp{\Om}$ inner product. The $L^2$-gradient of
$\Jt$ at $m$ is defined by
\begin{equation}
  \grad \Jt(m) := \Rmap_{L^2}(\mathrm{d}\Jt(m)) \in \Mspace,
\end{equation}
that is, $\grad \Jt(m)$ is the unique element of $\Mspace$ satisfying
\begin{equation}
  \mathrm{d}\Jt(m;\, \delta m)
  = \inner{\grad \Jt(m)}{\delta m}{\Lp{\Om}}
  \qquad \forall \delta m \in \Mspace.
\end{equation}
\end{defn}

Linearity of the inner product lets us factor
\cref{eq:dJtilde-pre-riesz} as
\begin{equation}
  \mathrm{d}\Jt(m;\, \delta m)
  = \inner{\lambda + \alpha m}{\delta m}{\Lp{\Om}}
  \qquad \forall \delta m \in \Mspace.
\end{equation}
Comparing with the definition of $\grad \Jt$ identifies
\begin{equation}
    \grad \Jt(m) = \lambda + \alpha m.
  \label{eq:l2gradient-poisson}
\end{equation}

The role played by the choice of inner product on $\Mspace$ in
\cref{defn:l2gradient} is the central concern of
\cref{sec:mesh-dependence}: replacing $\Rmap_{L^2}$ by the
sequence-space Riesz map $\Rmap_{\ell^2}$ produces a numerically
different gradient and, consequently, a mesh-dependent optimisation
algorithm.


% =============================================================================
\section{Numerical Implementation and Verification of the Poisson Problem}
\label{sec:poisson-numerics}
% =============================================================================

Having derived the weak form of the forward problem, the adjoint
equation and the $L^2$-gradient of the reduced functional, we now
implement the Poisson optimal control problem numerically. The
implementation serves two purposes: first, to verify that the finite
element discretisation of the forward problem behaves as expected;
second, to confirm that the adjoint-based gradient computation is
correct before extending the same approach to more complicated
problems in \cref{sec:burgers-numerics,sec:mesh-dependence}.

We carry out the numerical work in two stages. First, we implement
the forward Poisson problem in plain Python using a simple $P_1$
finite element method on a triangular mesh; the hand-written
implementation makes the main finite element steps explicit, namely
mesh construction, basis functions, element-wise assembly and the
imposition of homogeneous Dirichlet boundary conditions. Second, we
implement the same forward problem in Firedrake~\citep{firedrake2016}
together with the adjoint solve, the gradient computation and the
Taylor test. The two implementations share the same mathematical
content; the first makes the finite element method explicit, and the
second expresses it more compactly using a domain-specific language.

\subsection{Forward Poisson by hand}
\label{sec:poisson-numerics:byhand}

We implement the forward problem in plain Python using a continuous
piecewise linear finite element method. The domain $\Om = (0, 1)^2$
is triangulated using a structured mesh obtained by splitting each
square cell into two triangles, and the discrete space is the standard
$P_1$ Lagrange space.

For each element, the implementation assembles the local stiffness
matrix and the local load vector by mapping from the physical triangle
to a reference triangle and applying numerical quadrature. The local
contributions are then accumulated into a global sparse linear system.
Since the problem is subject to homogeneous Dirichlet boundary
conditions, the algebraic system is modified so that boundary degrees
of freedom are fixed to zero.

The hand-written implementation prioritises clarity over efficiency.
Writing the assembly explicitly makes the connection between the weak
form derived in \cref{sec:poisson-weak} and the finite element
matrices that appear in the computation transparent.

\subsection{Forward Poisson solve in Firedrake}
\label{sec:poisson-numerics:firedrake}

We then implement the same forward problem in
Firedrake~\citep{firedrake2016} using continuous Galerkin elements of
degree one. The variational forms are written directly from the weak
formulation \cref{eq:weak-forward-compact}, and the homogeneous
Dirichlet boundary condition is imposed strongly on $\dom$.

The Firedrake implementation is markedly more concise than the plain
Python version, but represents the same finite element problem.
Working in Firedrake also lets us express the forward solve, the
adjoint solve and the gradient computation in a form that stays close
to the mathematical notation introduced earlier in the work; this is
critical for the adjoint-based experiments of
\cref{sec:mesh-dependence}, which rely on the algorithmic
differentiation provided by the pyadjoint extension to
Firedrake~\citep{mitusch2019}.

\subsection{Manufactured solution and convergence study}
\label{sec:poisson-numerics:mms}

To test the correctness of the forward implementations, we use the
method of manufactured solutions, following the presentation in
Chapter~6 of \citet{hc2026}. We prescribe a smooth function
satisfying the homogeneous boundary conditions and choose the source
term accordingly. We define
\begin{equation}
  u_{\mathrm{exact}}(x, y) = \sin(\pi x) \sin(\pi y),
\end{equation}
which vanishes on $\dom$. The corresponding source term is
\begin{equation}
  m(x, y) = -\Delta u_{\mathrm{exact}}(x, y)
         = 2\pi^{2} \sin(\pi x) \sin(\pi y).
\end{equation}
With this choice, the continuous problem admits the known exact
solution $u = u_{\mathrm{exact}}$.

Let $u_h$ denote the finite element approximation on a mesh of size
$h$. For sufficiently smooth solutions, standard finite element
theory predicts that the $\Lp{\Om}$ error for piecewise linear
elements satisfies
\begin{equation}
  \|u - u_h\|_{\Lp{\Om}} = \mathcal{O}(h^{2}).
\end{equation}
We estimate the experimental order of convergence between two
successive mesh levels by
\begin{equation}
  q = \frac{\log(e_H / e_h)}{\log(h_H / h_h)},
\end{equation}
where $e_h = \|u - u_h\|_{\Lp{\Om}}$.

\paragraph{Hand-written results.}
The errors and rates observed on the plain Python $P_1$ implementation
appear in \cref{tab:poisson-byhand-rates}.

\begin{table}[H]
  \centering
  \begin{tabular}{r S[table-format=1.4] S[table-format=1.4e-2] S[table-format=1.4]}
    \toprule
    {$n_x$} & {$h$} & {$L^2$ error} & {rate} \\
    \midrule
      4 & 0.2500 & 7.5963e-02 & \multicolumn{1}{c}{-} \\
      8 & 0.1250 & 2.0409e-02 & 1.8961 \\
     16 & 0.0625 & 5.2009e-03 & 1.9724 \\
     32 & 0.0312 & 1.3066e-03 & 1.9930 \\
     64 & 0.0156 & 3.2704e-04 & 1.9982 \\
    128 & 0.0078 & 8.1786e-05 & 1.9996 \\
    256 & 0.0039 & 2.0448e-05 & 1.9999 \\
    512 & 0.0020 & 5.1121e-06 & 2.0000 \\
    \bottomrule
  \end{tabular}
  \caption{Manufactured-solution convergence test for the
  hand-written $P_1$ implementation of the forward Poisson problem.
  The observed rates approach the expected value of $2$.}
  \label{tab:poisson-byhand-rates}
\end{table}

\paragraph{Firedrake results.}
The same test on a sequence of uniformly refined Firedrake meshes
produced the rates of \cref{tab:poisson-firedrake-rates}.

\begin{table}[H]
  \centering
  \begin{tabular}{r S[table-format=1.4] S[table-format=1.4e-2] S[table-format=1.4]}
    \toprule
    {$N$} & {$h$} & {$L^2$ error} & {rate} \\
    \midrule
      4 & 0.2500 & 6.2103e-02 & \multicolumn{1}{c}{-} \\
      8 & 0.1250 & 1.8332e-02 & 1.7603 \\
     16 & 0.0625 & 4.7854e-03 & 1.9376 \\
     32 & 0.0312 & 1.2095e-03 & 1.9842 \\
     64 & 0.0156 & 3.0321e-04 & 1.9960 \\
    128 & 0.0078 & 7.5855e-05 & 1.9990 \\
    256 & 0.0039 & 1.8967e-05 & 1.9998 \\
    512 & 0.0020 & 4.7420e-06 & 1.9999 \\
    \bottomrule
  \end{tabular}
  \caption{Manufactured-solution convergence test for the Firedrake
  implementation of the forward Poisson problem. The observed rates
  approach $2$, giving confidence that the Firedrake implementation
  of the forward problem is correct.}
  \label{tab:poisson-firedrake-rates}
\end{table}

\subsection{Adjoint solve and gradient computation}
\label{sec:poisson-numerics:adjoint}

The derivation of \cref{sec:poisson-derivative} states that the
gradient of the reduced functional is computed by first solving the
adjoint equation and then forming the $L^2$-gradient
$\grad \Jt(m) = \lambda + \alpha m$. The Firedrake implementation
realises this in three stages:
\begin{enumerate}
  \item solve the forward Poisson problem to obtain the state $u$;
  \item solve the adjoint problem to obtain $\lambda$;
  \item compute the gradient as the function $\lambda + \alpha m$.
\end{enumerate}
The numerical implementation mirrors the theoretical derivation. In
particular, it exploits the fact that the adjoint equation lets us
compute the derivative of the reduced functional without solving a
separate tangent linear problem for every perturbation direction.

\subsection{Taylor test for gradient verification}
\label{sec:poisson-numerics:taylor}

After implementing the adjoint and gradient, we carry out a Taylor
test to verify that the computed gradient is correct. The Taylor test
is a standard practical check for adjoint
implementations~\citep{mitusch2019}.

Let $\delta m$ be a perturbation direction and let $h > 0$ be a small
parameter. The reduced functional satisfies the first-order Taylor
expansion
\begin{equation}
  \Jt(m + h \delta m)
  = \Jt(m) + h\, \mathrm{d}\Jt(m;\, \delta m) + \mathcal{O}(h^{2}).
\end{equation}
If the gradient implementation is correct, the first-order remainder
\begin{equation}
  R_1(h)
  := \left|
       \Jt(m + h \delta m) - \Jt(m) - h\, \mathrm{d}\Jt(m;\, \delta m)
     \right|
\end{equation}
decays as $\mathcal{O}(h^{2})$ as $h \to 0$.

Using a smooth perturbation direction, we obtained the remainder
sequence reported in \cref{tab:poisson-taylor}.

\begin{table}[H]
  \centering
  \begin{tabular}{S[table-format=1.3e-1] S[table-format=1.4e-2] S[table-format=1.4]}
    \toprule
    {$h$} & {first-order remainder} & {observed rate} \\
    \midrule
    1.000e-1 & 1.8923e-6 & {-}   \\
    5.000e-2 & 4.7307e-7 & 2.0000 \\
    2.500e-2 & 1.1827e-7 & 2.0000 \\
    1.250e-2 & 2.9567e-8 & 2.0000 \\
    6.250e-3 & 7.3917e-9 & 2.0000 \\
    \bottomrule
  \end{tabular}
  \caption{Taylor test for the gradient of the reduced Poisson
  functional. The observed second-order decay of the Taylor remainder
  is strong evidence that the adjoint solve and the gradient
  computation are implemented correctly.}
  \label{tab:poisson-taylor}
\end{table}


% =============================================================================
\section{The Forward Problem for the Scalar Viscous Burgers Equation}
\label{sec:burgers-weak}
% =============================================================================

We now consider a nonlinear model problem, the scalar viscous Burgers
equation. Following the notation established earlier, we write the
state variable as $u$ and the control as $m$. This is a scalar
simplification of the time-dependent viscous Burgers model considered
by \citet{shfp2017}, whose treatment is vector-valued. The principal
new features compared with the Poisson problem are nonlinearity in $u$
(through the convection term), time dependence (which we treat with
backward Euler in \cref{sec:burgers-numerics}), and a
backward-in-time adjoint solve.

\subsection{Strong formulation}
\label{sec:burgers-weak:strong}

Let $\Om \subset \R$ be a bounded spatial domain and let $T > 0$. The
scalar viscous Burgers equation reads
\begin{equation}
  \label{eq:burgers-strong}
  \frac{\partial u}{\partial t}
  + u \frac{\partial u}{\partial x}
  - \nu \frac{\partial^{2} u}{\partial x^{2}} = m
  \qquad \text{in } (0, T) \times \Om,
\end{equation}
subject to homogeneous Dirichlet boundary conditions
\begin{equation}
  u = 0 \qquad \text{on } (0, T) \times \dom,
\end{equation}
and initial condition
\begin{equation}
  u(0, x) = u_0(x) \qquad \text{for } x \in \Om.
\end{equation}
Here, $\nu > 0$ is the viscosity parameter.

\subsection{Weak formulation}
\label{sec:burgers-weak:weak}

For each fixed $t \in (0, T)$, we take the spatial state space to be
$V := H_0^1(\Om)$. Let $v \in V$ be a test function. Multiplying
\cref{eq:burgers-strong} by $v$ and integrating over $\Om$ gives
\begin{equation}
  \int_{\Om} \frac{\partial u}{\partial t}\, v \dx
  + \int_{\Om} u \frac{\partial u}{\partial x}\, v \dx
  - \nu \int_{\Om} \frac{\partial^{2} u}{\partial x^{2}}\, v \dx
  = \int_{\Om} m v \dx.
\end{equation}
Integrating the diffusion term by parts and using $v = 0$ on
$\dom$,
\begin{equation}
  \label{eq:burgers-weak}
  \int_{\Om} \frac{\partial u}{\partial t}\, v \dx
  + \int_{\Om} u \frac{\partial u}{\partial x}\, v \dx
  + \nu \int_{\Om} \frac{\partial u}{\partial x}\,
                   \frac{\partial v}{\partial x} \dx
  = \int_{\Om} m v \dx
  \qquad \forall v \in V,
\end{equation}
for almost every $t \in (0, T)$.


% =============================================================================
\section{Derivative of the Reduced Functional for the Burgers Problem}
\label{sec:burgers-derivative}
% =============================================================================

We extend the adjoint-based derivative analysis of
\cref{sec:poisson-derivative} to the scalar Burgers problem. As before,
the reduced functional is
\begin{equation}
  \Jt(m) := J(u(m), m),
\end{equation}
now with the time-integrated tracking objective
\begin{equation}
  \label{eq:burgers-objective}
  J(u, m)
  := \tfrac{1}{2} \|u - d\|_{L^{2}((0, T) \times \Om)}^{2}
   + \tfrac{\alpha}{2} \|m\|_{L^{2}((0, T) \times \Om)}^{2}.
\end{equation}
We take the control space to be
\begin{equation}
  \Mspace := L^{2}((0, T) \times \Om),
\end{equation}
and the test space to be
\begin{equation}
  \Vtest := L^{2}(0, T;\, V).
\end{equation}
The forward model operator $F$ is the residual of
\cref{eq:burgers-weak}, in direct analogy with
\cref{defn:forward-operator}.

\subsection{Tangent linear equation}
\label{sec:burgers-deriv:tlm}

Let $\delta m \in \Mspace$ be a control perturbation and let
$\delta u := \mathrm{d} u(m;\, \delta m)$ denote the corresponding
variation in the state. Differentiating $F(u(m), m;\, v) = 0$ with
respect to $m$ yields
\begin{equation}
  \partial_u F(u, m;\, \delta u, v) + \partial_m F(u, m;\, \delta m, v) = 0
  \qquad \forall v \in \Vtest.
\end{equation}

The derivative of the nonlinear convection term is
\begin{equation}
  (u + \varepsilon \delta u)
  \frac{\partial}{\partial x}(u + \varepsilon \delta u)
  = u \frac{\partial u}{\partial x}
  + \varepsilon
    \left(
      \delta u \frac{\partial u}{\partial x}
      + u \frac{\partial \delta u}{\partial x}
    \right)
  + \mathcal{O}(\varepsilon^{2}),
\end{equation}
so that
\begin{equation}
  \partial_u F(u, m;\, \delta u, v)
  = \int_0^T \! \int_{\Om}
    \left(
      \frac{\partial \delta u}{\partial t} v
      + \left( \delta u \frac{\partial u}{\partial x}
               + u \frac{\partial \delta u}{\partial x} \right) v
      + \nu \frac{\partial \delta u}{\partial x}
            \frac{\partial v}{\partial x}
    \right) \dx \dt,
\end{equation}
while
\begin{equation}
  \partial_m F(u, m;\, \delta m, v)
  = - \int_0^T \! \int_{\Om} \delta m\, v \dx \dt.
\end{equation}

Therefore, the tangent linear equation is
\begin{equation}
  \label{eq:burgers-tlm}
  \int_0^T \! \int_{\Om}
  \left(
    \frac{\partial \delta u}{\partial t} v
    + \left( \delta u \frac{\partial u}{\partial x}
             + u \frac{\partial \delta u}{\partial x} \right) v
    + \nu \frac{\partial \delta u}{\partial x}
          \frac{\partial v}{\partial x}
  \right) \dx \dt
  = \int_0^T \! \int_{\Om} \delta m\, v \dx \dt
  \qquad \forall v \in \Vtest,
\end{equation}
together with the initial condition $\delta u(0, x) = 0$.

\subsection{Adjoint equation}
\label{sec:burgers-deriv:adjoint}

The adjoint variable $\lambda$ is defined by
\begin{equation}
  \partial_u F(u, m;\, w, \lambda) = \partial_u J(u, m;\, w)
  \qquad \forall w \in \Vtest,
\end{equation}
where
\begin{equation}
  \partial_u J(u, m;\, w)
  = \int_0^T \! \int_{\Om} (u - d) w \dx \dt.
\end{equation}
Substituting the expression for $\partial_u F$,
\begin{equation}
  \int_0^T \! \int_{\Om}
  \left(
    \frac{\partial w}{\partial t} \lambda
    + \left( w \frac{\partial u}{\partial x}
             + u \frac{\partial w}{\partial x} \right) \lambda
    + \nu \frac{\partial w}{\partial x}
          \frac{\partial \lambda}{\partial x}
  \right) \dx \dt
  = \int_0^T \! \int_{\Om} (u - d) w \dx \dt
  \qquad \forall w \in \Vtest.
\end{equation}

We integrate the time-derivative term by parts in time,
\begin{equation}
  \int_0^T \! \int_{\Om} \frac{\partial w}{\partial t} \lambda \dx \dt
  = \int_{\Om} w(T, x) \lambda(T, x) \dx
  - \int_{\Om} w(0, x) \lambda(0, x) \dx
  - \int_0^T \! \int_{\Om} w \frac{\partial \lambda}{\partial t} \dx \dt.
\end{equation}
Admissible state variations satisfy $w(0, x) = 0$, so the initial
term vanishes. We impose the terminal condition
\begin{equation}
  \label{eq:burgers-terminal}
  \lambda(T, x) = 0.
\end{equation}

Integrating the convection term by parts in space,
\begin{equation}
  \int_{\Om}
  \left( w \frac{\partial u}{\partial x}
         + u \frac{\partial w}{\partial x} \right) \lambda \dx
  = - \int_{\Om} u \frac{\partial \lambda}{\partial x} w \dx,
\end{equation}
where the boundary term vanishes because of the homogeneous Dirichlet
condition. The weak adjoint equation therefore reads: find $\lambda$
such that
\begin{equation}
  \label{eq:burgers-adjoint-weak}
  \int_0^T \! \int_{\Om}
  \left(
    - \frac{\partial \lambda}{\partial t} w
    - u \frac{\partial \lambda}{\partial x} w
    + \nu \frac{\partial \lambda}{\partial x}
          \frac{\partial w}{\partial x}
  \right) \dx \dt
  = \int_0^T \! \int_{\Om} (u - d) w \dx \dt
  \qquad \forall w \in \Vtest,
\end{equation}
together with the terminal condition \cref{eq:burgers-terminal}.

The corresponding strong form is
\begin{equation}
  \label{eq:burgers-adjoint-strong}
  -\frac{\partial \lambda}{\partial t}
  - u \frac{\partial \lambda}{\partial x}
  - \nu \frac{\partial^{2} \lambda}{\partial x^{2}}
  = u - d
  \qquad \text{in } (0, T) \times \Om,
\end{equation}
with $\lambda = 0$ on $(0, T) \times \dom$ and
$\lambda(T, x) = 0$.

\subsection{Derivative of the reduced functional}
\label{sec:burgers-deriv:reduced}

Applying the chain rule as in the Poisson case,
\begin{equation}
  \mathrm{d}\Jt(m;\, \delta m)
  = \partial_u J(u, m;\, \delta u)
  + \partial_m J(u, m;\, \delta m),
\end{equation}
where
\begin{equation}
  \partial_m J(u, m;\, \delta m)
  = \alpha \int_0^T \! \int_{\Om} m\, \delta m \dx \dt.
\end{equation}
Using the adjoint equation,
\begin{equation}
  \partial_u J(u, m;\, \delta u)
  = \partial_u F(u, m;\, \delta u, \lambda),
\end{equation}
and the tangent linear equation with $v = \lambda$,
\begin{equation}
  \partial_u F(u, m;\, \delta u, \lambda)
  = -\partial_m F(u, m;\, \delta m, \lambda)
  = \int_0^T \! \int_{\Om} \delta m\, \lambda \dx \dt.
\end{equation}
Therefore
\begin{equation}
  \label{eq:burgers-dJtilde}
  \mathrm{d}\Jt(m;\, \delta m)
  = \int_0^T \! \int_{\Om} (\lambda + \alpha m)\, \delta m \dx \dt
  \qquad \forall \delta m \in \Mspace.
\end{equation}

\subsection{The \texorpdfstring{$L^2$}{L2}-gradient}
\label{sec:burgers-deriv:l2grad}

Since the control space is the Hilbert space
$\Mspace = L^{2}((0, T) \times \Om)$, the Riesz representation theorem
identifies the reduced derivative \cref{eq:burgers-dJtilde} with a
unique $L^2$-gradient. Comparing
\cref{eq:burgers-dJtilde} with
\begin{equation}
  \mathrm{d}\Jt(m;\, \delta m)
  = \inner{\grad \Jt(m)}{\delta m}{L^{2}((0, T) \times \Om)},
\end{equation}
we identify
\begin{equation}
    \grad \Jt(m) = \lambda + \alpha m.
\end{equation}
The reduced gradient has the same formal structure as in the Poisson
case (see \cref{eq:l2gradient-poisson}), but the adjoint variable
$\lambda$ is now determined by the backward-in-time adjoint equation
\cref{eq:burgers-adjoint-strong} which depends on the forward state
$u$.


% =============================================================================
\section{Numerical Implementation and Verification of the Burgers Problem}
\label{sec:burgers-numerics}
% =============================================================================

Having derived the weak formulation, the tangent linear and adjoint
equations and the $L^2$-gradient of the reduced functional for the
scalar viscous Burgers equation, we now implement the corresponding
optimal control problem numerically. The aim is the same as in
\cref{sec:poisson-numerics}: verifying that the discretisation of the
forward problem behaves as expected, and checking that the
adjoint-based gradient computation is correct.

In contrast to the Poisson chapter, where we provided both a plain
Python implementation and a Firedrake implementation, we treat the
Burgers problem at the lower level only. The main new difficulties
specific to Burgers' equation are the nonlinearity, the time dependence
and the backward-in-time adjoint solve. Implementing these ingredients
directly in plain Python makes the structure of the method more
transparent before moving to a higher-level framework.

\subsection{Scalar Burgers numerically}
\label{sec:burgers-numerics:byhand}

We consider the one-dimensional spatial domain $\Om = (0, 1)$ and a
final time $T > 0$. The implementation discretises the scalar viscous
Burgers equation in space using continuous piecewise linear finite
elements on a uniform mesh, and in time using the backward Euler
method.

Let $0 = t_0 < t_1 < \cdots < t_N = T$ denote the time grid, with
uniform step
\begin{equation}
  \Delta t = \frac{T}{N}.
\end{equation}
At each implicit time level, the nonlinear problem is solved using
Newton's method. If $U^n$ denotes the finite element coefficient
vector at time $t_n$, then the fully discrete forward problem at
$t_{n+1}$ has the form
\begin{equation}
  M \frac{U^{n+1} - U^{n}}{\Delta t}
  + C(U^{n+1}) + \nu K U^{n+1} = F^{n+1},
\end{equation}
where $M$ is the finite element mass matrix, $K$ is the stiffness
matrix, $C(U^{n+1})$ is the nonlinear convection contribution, and
$F^{n+1}$ is the load vector associated with the source term or
control at time $t_{n+1}$.

As with the plain Python Poisson implementation, this code prioritises
clarity over efficiency. Writing the assembly, the time stepping and
the Newton linearisation explicitly makes the path from the nonlinear
weak form to the algebraic system solved at each time step explicit.
The implementation imposes homogeneous Dirichlet boundary conditions
strongly at the two endpoints of the interval.

\subsection{Manufactured solution and convergence study}
\label{sec:burgers-numerics:mms}

To verify the forward implementation, we again use the method of
manufactured solutions. We prescribe a smooth exact solution
satisfying the homogeneous Dirichlet boundary conditions, and choose
the corresponding source term so that this exact solution satisfies
the Burgers equation.

We define
\begin{equation}
  u_{\mathrm{exact}}(t, x) = e^{-t} \sin(\pi x),
\end{equation}
which vanishes at $x = 0$ and $x = 1$ and therefore satisfies the
boundary conditions. Substituting into the strong form
\cref{eq:burgers-strong} gives the manufactured source term
\begin{equation}
  m_{\mathrm{exact}}(t, x)
  = - e^{-t} \sin(\pi x)
  + \pi e^{-2t} \sin(\pi x) \cos(\pi x)
  + \nu \pi^{2} e^{-t} \sin(\pi x).
\end{equation}

We solve the numerical problem on a sequence of uniformly refined
meshes with $N \in \{10, 20, 40, 80\}$ elements. Since the problem is
time-dependent, the temporal discretisation is refined at the same
rate by taking $N_t = 8N$, so that $\Delta t$ remains proportional to
the mesh size $h$.

We record two error measures: the final-time $\Lp{\Om}$ error,
\begin{equation}
  e_h^{\mathrm{final}} := \|u_h(T) - u_{\mathrm{exact}}(T)\|_{\Lp{\Om}},
\end{equation}
and a fully discrete space-time error,
\begin{equation}
  e_h^{\mathrm{st}}
  := \left( \Delta t \sum_{n=1}^{N_t}
              \|u_h^n - u_{\mathrm{exact}}(t_n)\|_{\Lp{\Om}}^{2}
        \right)^{1/2}.
\end{equation}
The experimental order of convergence between two successive
refinements is
\begin{equation}
  q = \frac{\log(e_H / e_h)}{\log(h_H / h_h)}.
\end{equation}

The observed results appear in \cref{tab:burgers-mms}.

\begin{table}[H]
  \centering
  \begin{tabular}{r r r S[table-format=1.3e-1] S[table-format=1.4]
                            S[table-format=1.3e-1] S[table-format=1.4]}
    \toprule
    {$N$} & {$N_t$} & {$h$}
      & {$e_h^{\mathrm{final}}$} & {rate}
      & {$e_h^{\mathrm{st}}$}    & {rate} \\
    \midrule
    10 &  80 & 0.1000 & 5.840e-3 & {-}   & 1.838e-3 & {-}   \\
    20 & 160 & 0.0500 & 1.452e-3 & 2.0077 & 4.590e-4 & 2.0027 \\
    40 & 320 & 0.0250 & 3.577e-4 & 2.0197 & 1.140e-4 & 2.0098 \\
    80 & 640 & 0.0125 & 8.679e-5 & 2.0414 & 2.810e-5 & 2.0209 \\
    \bottomrule
  \end{tabular}
  \caption{Manufactured-solution convergence test for the scalar
  Burgers forward solver. Both error measures decrease at an observed
  rate very close to $2$, consistent with the expected behaviour of
  the fully discrete scheme for this smooth manufactured solution.}
  \label{tab:burgers-mms}
\end{table}

\subsection{Discrete adjoint solve and gradient computation}
\label{sec:burgers-numerics:adjoint}

We implement the discrete adjoint equation and compute the gradient
of the reduced functional at the fully discrete level, consistent
with the backward Euler time discretisation of the forward problem.

If $D^n$ denotes the discrete desired state at time $t_n$ and $m^n$
the discrete control, the fully discrete tracking functional is
\begin{equation}
  J_h(U, m)
  = \frac{\Delta t}{2} \sum_{n=1}^{N_t} \|U^n - D^n\|_M^{2}
  + \frac{\alpha \Delta t}{2} \sum_{n=1}^{N_t} \|m^n\|_M^{2},
\end{equation}
where $\|v\|_M^{2} := v^{T} M v$ is the mass-matrix representation
of the $\Lp{\Om}$ norm.

We generate the desired state by first choosing a smooth reference
control, solving the forward Burgers problem for that control, and
taking the resulting discrete state trajectory as the target
trajectory. We then choose a different trial control and solve the
forward problem again in order to compute the state, the objective
value, the adjoint and the gradient.

As expected for a time-dependent problem, the discrete adjoint
variables are computed by solving a sequence of linear systems
backward in time, starting from a zero terminal condition. The
forward Burgers problem is nonlinear and is solved forward in time,
while the adjoint problem is linearised about the computed state and
is solved backward in time \citep[Sect.~1.5-1.6]{shfp2017}.

The resulting discrete gradient has the form
\begin{equation}
  g^n = \alpha m^n + \frac{1}{\Delta t} p^n,
\end{equation}
where $p^n$ is the discrete adjoint variable at time level $t_n$.
This is the fully discrete analogue of the continuous expression
$\grad \Jt(m) = \lambda + \alpha m$ derived earlier.

For the representative trial control used in the experiments, the
computed objective value was
\begin{equation}
  J_h(m) = 1.323574 \times 10^{-4},
\end{equation}
and the average number of Newton iterations required in the forward
solve was approximately $3$ per time step, indicating that the
nonlinear solver converged reliably for the chosen problem parameters.

\subsection{Taylor test for gradient verification}
\label{sec:burgers-numerics:taylor}

After implementing the discrete adjoint and gradient, we carry out a
Taylor test in the same spirit as in
\cref{sec:poisson-numerics:taylor}.

Let $\delta m$ be a perturbation direction in the discrete control
space and let $h > 0$. The reduced functional satisfies
\begin{equation}
  \Jt_h(m + h \delta m)
  = \Jt_h(m) + h\, \mathrm{d}\Jt_h(m;\, \delta m)
  + \mathcal{O}(h^{2}),
\end{equation}
and if the gradient implementation is correct, the second-order
remainder
\begin{equation}
  R_2(h)
  := \left|
       \Jt_h(m + h \delta m) - \Jt_h(m)
       - h\, \mathrm{d}\Jt_h(m;\, \delta m)
     \right|
\end{equation}
decays as $\mathcal{O}(h^{2})$ as $h \to 0$.

Using a smooth perturbation direction normalised in the discrete
control inner product, we obtained the results of
\cref{tab:burgers-taylor}.

\begin{table}[H]
  \centering
  \begin{tabular}{S[table-format=1.3e-1] S[table-format=1.4e-1] S[table-format=1.4]}
    \toprule
    {$h$} & {second-order remainder} & {observed rate} \\
    \midrule
    1.000e-1 & 6.7317e-5 & {-}   \\
    5.000e-2 & 1.6828e-5 & 2.0001 \\
    2.500e-2 & 4.2068e-6 & 2.0001 \\
    1.250e-2 & 1.0517e-6 & 2.0000 \\
    6.250e-3 & 2.6291e-7 & 2.0000 \\
    \bottomrule
  \end{tabular}
  \caption{Taylor test for the discrete reduced functional in the
  scalar Burgers problem. The observed second-order decay of the
  Taylor remainder is strong evidence that the discrete adjoint solve
  and gradient computation are implemented correctly.}
  \label{tab:burgers-taylor}
\end{table}

\subsection{Summary}
\label{sec:burgers-numerics:summary}

This section extended the numerical verification pipeline from the
Poisson problem to a nonlinear, time-dependent Burgers problem. We
implemented the forward scalar viscous Burgers equation in plain
Python using $P_1$ finite elements in space, backward Euler in time
and Newton's method for the nonlinear solve. A manufactured-solution
convergence study showed second-order decay in both the final-time
$L^{2}$ error and a discrete space-time error. We then implemented the
discrete adjoint equation and the reduced gradient, and a Taylor test
confirmed that the gradient computation is correct.

Together, the verifications of
\cref{sec:poisson-numerics,sec:burgers-numerics} show that the
adjoint-based optimisation framework developed in
\cref{sec:poisson-derivative,sec:burgers-derivative} extends
successfully from a linear, stationary problem to a nonlinear,
time-dependent one. This provides a bridge between the simple Poisson
example and the broader questions of PDE-constrained optimisation and
mesh dependence treated in the reference text \citep{shfp2017} and
taken up next in \cref{sec:mesh-dependence}.


% =============================================================================
\section{Mesh Dependence in PDE-Constrained Optimisation}
\label{sec:mesh-dependence}
% =============================================================================

So far, we have derived, discretised and verified the
adjoint-based gradient for both the Poisson and the scalar viscous
Burgers optimal control problems. In each case the $L^2$-gradient
$\grad \Jt(m) = \lambda + \alpha m$ was identified via the Riesz
representation theorem (\cref{thm:riesz}). That identification is not
unique: it depends on the choice of Hilbert space structure placed on
the control space $\Mspace$. In this section we examine the practical
consequences of that choice when the reduced functional is minimised
numerically. Following \citet[Chap.~2]{shfp2017}, we show that two
different but mathematically valid choices of inner product yield
optimisation algorithms whose iteration counts behave very differently
as the mesh is refined.

\subsection{Background and motivation}
\label{sec:meshdep:background}

\Cref{sec:prelim:riesz,sec:deriv:l2grad} established that the
Fr{\'e}chet derivative $\mathrm{d}\Jt(m)$ is an element of the dual
space $\Mspace^{*}$, and that the Riesz map
$\Rmap_H : \Mspace^{*} \to \Mspace$ identifies each derivative with a
primal element of $\Mspace$. When $\Mspace = \Lp{\Om}$ is treated as
an $L^{2}$ Hilbert space, the Riesz map $\Rmap_{L^2}$ corresponds in
the discrete setting to multiplication by the inverse of the Galerkin
mass matrix \citep[Eq.~(1.125)]{shfp2017}. When the same coefficient
vector is regarded instead as an element of $\R^{d}$ under the
standard Euclidean ($\ell^{2}$) inner product, no mass matrix appears
and the Riesz map is the identity.

The two choices yield primal gradients that differ at every mesh
refinement: the $\ell^{2}$ representer carries the local cell sizes
encoded in the diagonal of the mass matrix, while the $L^{2}$
representer does not. Search directions and stopping criteria based
on gradient norms therefore behave differently in the two settings.
\Citet[Sect.~2.1]{shfp2017} note that ``many well-established
continuous optimisation packages are not designed for function-based
optimisation and are hard-coded to apply the Euclidean ($\ell^{2}$)
inner product'', and explicitly list \texttt{scipy.optimize}, IPOPT
and TAO among them. The Moola library and recent versions of
pyadjoint allow the user to specify the Hilbert space inner product
on the control space.

\subsection{Hypothesis}
\label{sec:meshdep:hypothesis}

We test the prediction of \citet[Sect.~2.2]{shfp2017}: for a fixed
gradient-norm convergence threshold, the iteration count of a
quasi-Newton method depends on which Riesz map is used to interpret
the derivative as a primal-space gradient. Concretely:

\begin{description}
  \item[H1.] When the $\ell^{2}$ Riesz map is used, the iteration
        count grows with the mesh ratio $h_{\max} / h_{\min}$
        (mesh-dependent convergence).
  \item[H2.] When the $L^{2}$ Riesz map is used, the iteration count
        is bounded uniformly in $h_{\max} / h_{\min}$
        (mesh-independent convergence).
\end{description}

\subsection{Mesh generation}
\label{sec:meshdep:meshes}

Non-uniform meshes of $\Om = (0, 1)^{2}$ with controlled
$h_{\max}/h_{\min}$ ratios are constructed using the open-source
mesher Netgen~\citep{schoeberl1997}, following the same recipe as
\citet[Sect.~2.3]{shfp2017}. The unit square is partitioned into a
coarse outer region with target cell size
$h_{\max} = 0.2$ and a finer inner sub-square
$(0, 0.4)^{2}$ with target cell size $h_{\max}/r$, where $r$ is the
prescribed ratio. The realised ratio is read back from the resulting
cell-diameter field of the assembled mesh and is generally somewhat
larger than the prescribed $r$, since Netgen smoothly grades the
transition between the two regions.

The function \texttt{graded\_unit\_square} in
\texttt{meshdep.meshes} encapsulates the recipe. The Netgen Python
package does not provide aarch64 binary distributions at the time of
writing and the source distribution does not build on the machine
used for the experiments. We therefore generate the six meshes for
$r \in \{4, 8, 16, 32, 64, 128\}$ once, on an x86\_64 machine, and
store them as Gmsh-format \texttt{.msh} files in the
\texttt{mesh\_generation/} directory of the project. The experiment
scripts then load the meshes via
\texttt{graded\_unit\_square\_from\_file}. The Gmsh header version
emitted by Netgen is \texttt{2.0}, which the Firedrake mesh reader
rejects with the message \texttt{``Gmsh file version 2.0 must be at
least 2.2''}; the binary layout of the two versions is identical, so
each file's header is patched from \texttt{2.0} to \texttt{2.2}
before use. The numbers of mesh elements and the realised ratios are
reported in \cref{tab:meshdep:mesh_summary}.

\begin{table}[H]
  \centering
  \begin{tabular}{r r S[table-format=3.2]}
    \toprule
    {target $r$} & {\# cells} & {realised $h_{\max}/h_{\min}$} \\
    \midrule
      4 &     244 &   5.92 \\
      8 &     771 &   9.99 \\
     16 &   2\,664 &  23.50 \\
     32 &   9\,764 &  41.89 \\
     64 &  36\,450 & 115.17 \\
    128 & 141\,138 & 193.88 \\
    \bottomrule
  \end{tabular}
  \caption{The six Netgen graded meshes used in the experiments.
  Realised ratios are measured from the cell-diameter field via
  \texttt{CellDiameter} in Firedrake.}
  \label{tab:meshdep:mesh_summary}
\end{table}

\subsection{Implementation}
\label{sec:meshdep:implementation}

The numerical experiments use the Python package \texttt{meshdep}
developed for this dissertation, built on top of
Firedrake~\citep{firedrake2016} and pyadjoint~\citep{mitusch2019}.
The Poisson optimal control problem of \cref{sec:poisson-derivative}
is exposed through
\texttt{meshdep.problems.poisson\_control.solve\_poisson\_control},
which accepts a \texttt{riesz\_map} argument with admissible values
\texttt{"l2"}, \texttt{"L2"} and \texttt{"H1"}. Internally the
argument is passed to pyadjoint's \texttt{Control} constructor:

\begin{lstlisting}[language=Python]
Jhat = ReducedFunctional(
    J, Control(m, riesz_map=riesz_map), tape=tape,
)
\end{lstlisting}

The keyword chooses which discrete Riesz map pyadjoint applies when
it converts the dual-space derivative $\mathrm{d}\Jt(m)$ to a primal
gradient. The forward Poisson solve uses a direct LU solver, so that
the discrete gradient is computed to machine precision and the
convergence threshold $\varepsilon = 10^{-7}$ on the gradient norm
$\|\Rmap_H(J')\|_H$ is not blocked by Krylov-solver noise.

Three points of implementation merit explicit discussion, because
each was non-trivial.

\paragraph{The convergence criterion is panel-wide, not
algorithm-specific.} \Citet[p.~67-68]{shfp2017} stop every algorithm
in a given panel on the same gradient norm
$\|\Rmap_H(J')\|_H \le 10^{-7}$, irrespective of the inner product
used internally by the algorithm. For an algorithm that does not
natively support a custom inner product (such as SciPy's L-BFGS-B),
\citet{shfp2017} note that ``we compute the functional derivative
with respect to the appropriate inner product outside the
$\ell^{2}$-employing algorithms''. We reproduce this externally
checked convergence criterion in
\texttt{meshdep.optimisers.solve\_with\_scipy\_external\_check}, which
runs SciPy's L-BFGS-B with its native tolerances set to $10^{-30}$
(so that they never trigger) and uses a callback that re-evaluates
$\Rmap_H(J')$ at the current iterate, computes its $H$-norm and
raises a sentinel exception once the threshold is met.

\paragraph{Routing through PETSc TAO/LMVM is possible, but requires
two non-default settings.} The natural attempt to reproduce the
Moola row of Table~2.2 in \citet{shfp2017} is to construct
\texttt{Control(m, riesz\_map="L2")} and hand the resulting
\texttt{ReducedFunctional} to pyadjoint's TAO interface. Pyadjoint's
TAO wrapper then installs the Riesz inverse $M_h^{-1}$ both as the
gradient-norm metric used in the convergence test (via
\texttt{tao.setGradientNorm}) and as the initial Hessian
approximation $B_{0}$ in the L-BFGS recursion (via
\texttt{setLMVMH0}). The LMVM matrix object then maintains the
limited-memory secant pairs in the inner product induced by $B_0$,
so the algorithm is, in principle, doing function-space L-BFGS.

In practice, two implementation issues need to be addressed before
this combination is competitive with the
Moola row of \citet[Table~2.2(a)]{shfp2017}. The first is purely
mechanical: pyadjoint's TAO wrapper accepts a Python \texttt{parameters}
dictionary, but the entries of that dictionary are not propagated to
the LMVM sub-Mat that PETSc allocates internally. The option
\texttt{tao\_lmvm\_mat\_lmvm\_hist\_size}, which controls the number
of secant pairs that the L-BFGS update keeps in memory, is silently
ignored when passed through the dictionary. It must be registered
on the global PETSc options database \emph{before} the
\texttt{TAOSolver} is instantiated:

\begin{lstlisting}[language=Python]
PETSc.Options().setValue('-tao_lmvm_mat_lmvm_hist_size', 100)
\end{lstlisting}

The second issue is that the PETSc default history size of five is
too small for the discrete Hessian on the strongly graded Netgen
meshes of \cref{tab:meshdep:mesh_summary}. A rank-five quasi-Newton
update can only reproduce a Hessian whose top five eigenvalues
already cover most of the spectrum; on the graded meshes the
discrete Hessian has a much longer tail, and the LMVM iteration
proceeds with a poorly conditioned search direction. The
consequence is slow convergence rather than outright failure - with
enough outer iterations, the gradient norm does reach machine
precision (we verified this in
\cref{sec:meshdep:tao-revisited}) - but the iteration count at
the panel-wide threshold $\varepsilon = 10^{-7}$ grows from roughly
$130$ on a Cartesian mesh to several hundred at higher gradings.
Setting the history to one hundred recovers a mesh-independent
iteration count of around thirty across the entire range.

\paragraph{A transparent Hilbert-space L-BFGS reference
implementation.} TAO/LMVM with the corrections above reproduces
Schwedes' qualitative pattern, but the rank-two update lives inside
a PETSc C kernel and the precise inner products used in the
two-loop recursion are not easy to inspect from Python. To make the
role of the inner product completely explicit, and to obtain a
reference implementation that we can read line by line, we
additionally implement BFGS in a Hilbert space directly. Following
\citet[Sect.~1.5.4.1]{shfp2017}, the function
\texttt{meshdep.optimisers.solve\_with\_hilbert\_lbfgs} implements an
L-BFGS algorithm in which every dot product used inside the
two-loop recursion uses the chosen Hilbert space inner product
$\inner{\cdot}{\cdot}{H}$ rather than the ambient $\ell^{2}$ inner
product on coefficient vectors. This is the distinguishing feature
of Schwedes' ``Moola'' row, and the algorithm is presented in
\cref{sec:meshdep:hilbert-lbfgs}.

\subsection{BFGS in a Hilbert space}
\label{sec:meshdep:hilbert-lbfgs}

Let $H$ denote either $\Lp{\Om}$ or $H^{1}(\Om)$, with inner product
$\inner{\cdot}{\cdot}{H}$ given respectively by
\begin{equation}
  \inner{a}{b}{L^{2}} = \int_{\Om} a\, b \dx,
  \qquad
  \inner{a}{b}{H^{1}}
     = \int_{\Om} \bigl(a\, b + \grad a \cdot \grad b\bigr) \dx.
\end{equation}
Given a current iterate $m_{k} \in \Mspace$ and the primal Riesz
representer $g_{k} := \Rmap_H(\mathrm{d}\Jt(m_{k}))$ of the
derivative at $m_{k}$, an L-BFGS step in $H$ is computed by the
following two-loop recursion. We store the most recent
$\ell \le 5$ curvature pairs
\begin{equation}
  s_{j} := m_{j+1} - m_{j} \in \Mspace,
  \qquad
  y_{j} := g_{j+1} - g_{j} \in \Mspace,
  \qquad
  \rho_{j} := \frac{1}{\inner{s_{j}}{y_{j}}{H}}.
\end{equation}

The search direction $p_{k}$ is computed by
\begin{align}
  q & \leftarrow g_{k}, \\
  \intertext{for $j = k-1, k-2, \ldots, k-\ell$:}
  \alpha_{j} & \leftarrow \rho_{j}\, \inner{s_{j}}{q}{H}, \\
  q          & \leftarrow q - \alpha_{j}\, y_{j}, \\
  \intertext{then}
  z          & \leftarrow \gamma_{k}\, q, \\
  \intertext{for $j = k-\ell, \ldots, k-1$:}
  \beta_{j}  & \leftarrow \rho_{j}\, \inner{y_{j}}{z}{H}, \\
  z          & \leftarrow z + (\alpha_{j} - \beta_{j})\, s_{j}, \\
  \intertext{and finally}
  p_{k}      & \leftarrow -z.
\end{align}
The initial Hessian approximation is rescaled at each iteration by
\begin{equation}
  \label{eq:meshdep:gamma_k}
  \gamma_{k}
  = \frac{\inner{s_{k-1}}{y_{k-1}}{H}}{\inner{y_{k-1}}{y_{k-1}}{H}},
\end{equation}
the standard L-BFGS rescaling \citep[Sect.~1.5.4.2]{shfp2017} which
adapts $B_{0}$ to local curvature. Empirically, omitting this step
causes the algorithm to plateau several orders of magnitude above
the convergence threshold, even on a well-conditioned mesh.

The line search uses a backtracking Armijo rule: starting from
$\alpha = 1$, we halve $\alpha$ until the sufficient-decrease
condition
\begin{equation}
  \Jt(m_{k} + \alpha\, p_{k})
    \le \Jt(m_{k})
       + c_{1}\, \alpha\, \inner{g_{k}}{p_{k}}{H},
  \qquad c_{1} = 10^{-4},
\end{equation}
is satisfied. Curvature pairs that violate the positive-curvature
condition $\inner{s_{j}}{y_{j}}{H} > 0$ are silently dropped from
the history. Convergence is declared once
$\|g_{k}\|_{H} \le \varepsilon = 10^{-7}$.

The Riesz solve used to obtain $g_{k}$ from the dual derivative
$\mathrm{d}\Jt(m_{k})$ is implemented internally with a direct LU
factorisation of the discrete inner-product matrix
($M_{h}$ for $H = L^{2}$, $M_{h} + K_{h}$ for $H = H^{1}$). The
default Firedrake \texttt{RieszMap} uses a Krylov solver, which is
adequate for forward Riesz representations but introduces an extra
solver tolerance into the convergence test of the optimisation
algorithm. Using a direct factorisation removes that tolerance from
the analysis entirely, so that any failure to reach the threshold
$\varepsilon = 10^{-7}$ is attributable to the optimisation
algorithm rather than to the linear solver used in the
gradient-norm computation.

\subsection{Numerical experiments: Panel (a), \texorpdfstring{$H = \Lp{\Om}$}{H = L2}}
\label{sec:meshdep:panel-a}

We now reproduce panel (a) of Table~2.2 in \citet{shfp2017}.
The state equation is the weak Poisson problem of
\cref{eq:weak-forward-compact}, the objective functional is
\cref{eq:objective-J} with $d(x, y) = \sin(\pi x) \sin(\pi y)$ and
$\alpha = 10^{-4}$, the control is initialised at $m \equiv 0$, and
the convergence threshold is $\varepsilon = 10^{-7}$ on
$\|\Rmap_{L^2}(J')\|_{L^{2}}$. Two algorithms are run, both with an
L-BFGS history of five:
\begin{itemize}
  \item SciPy's L-BFGS-B run on the $\ell^{2}$ coefficient vector
        through
        \texttt{meshdep.optimisers.solve\_with\_scipy\_external\_check},
        with the convergence criterion checked externally in
        $L^{2}$;
  \item the Hilbert-space L-BFGS of
        \cref{sec:meshdep:hilbert-lbfgs}, run with
        $H = \Lp{\Om}$ through
        \texttt{meshdep.optimisers.solve\_with\_hilbert\_lbfgs}.
\end{itemize}

The iteration counts at convergence are reported in
\cref{tab:meshdep:panel-a}. We observed clear mesh-dependent growth
of the $\ell^{2}$ row and a flat $L^{2}$ row, in line with
hypotheses H1 and H2 of \cref{sec:meshdep:hypothesis}.

\begin{table}[H]
  \centering
  \begin{tabular}{l l r r r r r}
    \toprule
    inner product & implementation
       & \multicolumn{5}{c}{target $h_{\max} / h_{\min}$} \\
    \cmidrule(lr){3-7}
       &        & 4 & 8 & 16 & 32 & 64 \\
    \midrule
    $\ell^{2}$ & SciPy (L-BFGS-B, external check)
                                 & 46 & 96 & 195 & 377 & 786 \\
    $L^{2}$    & Hilbert-space L-BFGS
                                 &  8 &  8 &   7 &   8 &   7 \\
    \bottomrule
  \end{tabular}
  \caption{Iteration counts at convergence for the Poisson optimal
  control problem on the Netgen graded meshes of
  \cref{tab:meshdep:mesh_summary}. Both rows used the same problem
  data, the same forward LU solver, the same convergence threshold
  $\varepsilon = 10^{-7}$ on $\|\Rmap_{L^2}(J')\|_{L^{2}}$ and the
  same L-BFGS history of five. The only difference between the rows
  is the inner product used internally by the algorithm.}
  \label{tab:meshdep:panel-a}
\end{table}

The $\ell^{2}$ row very nearly doubles at every refinement
($\times 2.09, \times 2.03, \times 1.93, \times 2.08$), which is
consistent with the polynomial-in-$h_{\max}/h_{\min}$ growth derived
analytically in \citet[Sect.~2.2.4]{shfp2017}. The $L^{2}$ row
fluctuates only between seven and eight iterations across a
twenty-fold change in realised mesh ratio. For reference,
\citet[Table~2.2(a)]{shfp2017} reports $27, 47, 75, 97, 133$ for the
SciPy/L-BFGS-B row and $22, 20, 22, 23, 23$ for the Moola/BFGS row at
the same target ratios. The qualitative pattern reproduces (the
$\ell^{2}$ row grows, the function-space row is flat); the absolute
counts differ from \citet{shfp2017} because the realised
$h_{\max}/h_{\min}$ on our meshes is somewhat larger than the
nominal target (cf.~\cref{tab:meshdep:mesh_summary}) and because the
backtracking Armijo line search converges in fewer outer iterations
than the strict-Wolfe line search used by Moola.

We separately ran the Hilbert-space L-BFGS at the largest mesh ratio
$r = 128$ (realised $193.88$, with $141\,138$ cells); the algorithm
converged in nine iterations, confirming that the flat-line
behaviour extends beyond the range reported by \citet{shfp2017}.

\subsection{Numerical experiments: Panel (b), \texorpdfstring{$H = H^{1}(\Om)$}{H = H1}}
\label{sec:meshdep:panel-b}

The same experiment is repeated under the $H^{1}$ inner product for
both the algorithm and the convergence criterion. \Citet[Table~2.2(b)]{shfp2017}
reports that the corresponding Moola row converges
in exactly four iterations across the entire range of mesh ratios.
We ran the Hilbert-space L-BFGS of \cref{sec:meshdep:hilbert-lbfgs}
with $H = H^{1}(\Om)$ on the Netgen meshes of
\cref{tab:meshdep:mesh_summary} for $r \in \{4, 8\}$ and obtained
iteration counts of $213$ and $286$, respectively. The algorithm
reached the same minimum value of the discrete objective as the
SciPy row (to six significant figures), but required substantially
more iterations than the four reported by \citet{shfp2017} for the
corresponding Moola row.

Two factors compound here. First, the panel-wide threshold is
imposed on $\|\Rmap_H(J')\|_H$, and because the embedding
$H^{1}(\Om) \hookrightarrow L^{2}(\Om)$ implies
$\|\cdot\|_{H^{1}} \ge \|\cdot\|_{L^{2}}$, the same numerical value
$\varepsilon = 10^{-7}$ is a strictly more demanding criterion in
$H^{1}$ than in $L^{2}$. Second, the Hessian of $\Jt$ in the $H^{1}$
inner product is poorly conditioned. The exact-arithmetic
eigenvalues of $(I - \Delta)^{-1}$ on the unit square scale as
$1/(1 + \pi^{2}(j^{2} + k^{2}))$, and composing this with the
$L^{2}$-Hessian of $\Jt$ produces a transformed Hessian whose
eigenvalues decay as $k^{-6}$ on the high-frequency eigenfunctions.
A limited-memory BFGS with history five therefore cannot resolve
the high-frequency modes in a constant number of iterations. The
iteration counts $213$ and $286$ are not a defect of the
implementation: convergence is monotone in the discrete objective,
and the algorithm terminates at the same minimiser as the SciPy
row.

That the bottleneck is the convergence criterion in $H^{1}$ rather
than $H^{1}$ \emph{per se} is confirmed by routing the same
\texttt{Control(m, riesz\_map="H1")} object through PETSc TAO/LMVM
instead. With the inner product matched to the threshold (the
gradient is checked in $H^{1}$ and the tolerance is loosened to
$10^{-3}$, the value for which $\|\Rmap_{H^{1}}(J')\|_{H^{1}}
\approx 10^{-3}$ corresponds to roughly the same primal accuracy as
$\|\Rmap_{L^{2}}(J')\|_{L^{2}} \approx 10^{-7}$ at the discrete
minimum), TAO/LMVM converges in $26$, $30$ and $23$ iterations at
realised mesh ratios $1$, $6.9$ and $18.3$. So the H$^{1}$ Riesz
map does deliver mesh-independent behaviour when paired with a
threshold expressed in the same inner product; the apparent failure
above is a mismatch between the inner product of the algorithm and
the inner product of the stopping criterion.

The four-iteration count reported by \citet{shfp2017} for the
$H^{1}$ Moola row is not explained in the book, and we have not
identified a Moola-specific preconditioning step that would
reproduce it from the algorithmic ingredients described in
\citet[Sect.~1.5.4.1]{shfp2017}. We conjecture either that Moola
applies an additional preconditioner that is not documented in the
book (such as solving the compact term of the $H^{1}$-Hessian
explicitly), or that its iteration count is defined relative to a
different convergence criterion than the one stated in the table
caption. Resolving the discrepancy would require inspection of the
Moola source code, which currently cannot be installed in the
Firedrake environment used for these experiments because the
distribution hard-codes a \texttt{dolfin} import at module-load
time.

\subsection{TAO/LMVM revisited: diagnosing the original stagnation}
\label{sec:meshdep:tao-revisited}

The discussion in \cref{sec:meshdep:implementation} reported that
routing the Poisson control problem through pyadjoint's TAO/LMVM
interface with \texttt{Control(m, riesz\_map="L2")} fails to
reproduce the Moola row, and that the failure can be removed by
enlarging the LMVM history. This subsection records the diagnostic
work that led to this conclusion, because the original failure mode
was subtle enough to be misdiagnosed once already.

Our first hypothesis was that the LMVM iteration was stagnating at
a \emph{noise floor} - a gradient norm below which round-off in the
adjoint tape prevents further reduction. The hypothesis was tested
by replacing TAO/LMVM with TAO/NLS (a Newton line search using
exact CG-based Hessian-vector products) and rerunning the same
graded sweep with an artificially tight tolerance
$\varepsilon = 10^{-30}$. On a quadratic problem one Newton step
lands at the minimum, so the iterate produced by NLS at termination
is the discrete optimum to within the tolerance of the Newton
linear solver. Evaluating the gradient at that iterate via the
adjoint tape gave $\|\Rmap_{L^{2}}(J')\|_{L^{2}} \approx 10^{-18}$
on every mesh in the sweep. The adjoint-tape noise floor is
therefore at machine precision, not at the $10^{-5}$ at which
LMVM stagnated. The original stagnation was not a noise floor.

The second hypothesis was that the limited-memory approximation
itself was responsible. To test this we ran LMVM with the same
artificially tight tolerance $\varepsilon = 10^{-30}$ and a
budget of $3000$ iterations, sweeping the history size
$h \in \{5, 25, 50, 100, 200\}$. The final gradient norm
$\|\Rmap_{L^{2}}(J')\|_{L^{2}}$ at termination is summarised in
\cref{tab:meshdep:tao-history}. With the default history $h = 5$,
the iteration does reach $10^{-13}$ on the uniform mesh given a
sufficient budget, but stagnates eleven orders of magnitude
higher than that on a graded mesh of stretch $5$. Increasing the
history to $25$ already brings both meshes down to machine
precision, and increasing it further changes the final gradient
norm only marginally. The original failure was therefore caused
by a rank-five quasi-Newton approximation being too coarse to
match the spectrum of the discrete Hessian on the graded mesh;
with enough memory, LMVM does reach the discrete minimum to
machine precision on every mesh tested.

\begin{table}[H]
  \centering
  \begin{tabular}{c r r r r r}
    \toprule
    stretch
       & \multicolumn{5}{c}{LMVM history size $h$} \\
    \cmidrule(lr){2-6}
       & 5 & 25 & 50 & 100 & 200 \\
    \midrule
    $0$ & $9.4\!\times\!10^{-14}$
        & $6.4\!\times\!10^{-18}$
        & $4.6\!\times\!10^{-18}$
        & $2.7\!\times\!10^{-18}$
        & $1.8\!\times\!10^{-18}$ \\
    $5$ & $1.6\!\times\!10^{-11}$
        & $5.0\!\times\!10^{-18}$
        & $3.0\!\times\!10^{-18}$
        & $1.6\!\times\!10^{-18}$
        & $9.1\!\times\!10^{-19}$ \\
    \bottomrule
  \end{tabular}
  \caption{Final $\|\Rmap_{L^{2}}(J')\|_{L^{2}}$ produced by
  TAO/LMVM with \texttt{Control(m, riesz\_map="L2")} on the Poisson
  control problem, with the tolerance set to $10^{-30}$ and a
  budget of $3000$ iterations. Stretch $0$ is a uniform Cartesian
  mesh ($h_{\max}/h_{\min} = 1$) and stretch $5$ produces a graded
  mesh with $h_{\max}/h_{\min} \approx 30$. The default
  PETSc history size of $h = 5$ stagnates eleven orders of
  magnitude above the discrete optimum on the graded mesh;
  history sizes from $h = 25$ upwards reach machine precision on
  both.}
  \label{tab:meshdep:tao-history}
\end{table}

\paragraph{Pyadjoint options-database propagation.} Setting
\texttt{tao\_lmvm\_mat\_lmvm\_hist\_size} via the Python
\texttt{parameters} dictionary that pyadjoint's TAO wrapper accepts
does not, in our pyadjoint version, propagate down to the
LMVM sub-Mat: the options-prefix that pyadjoint installs does not
match the prefix that PETSc uses internally for the matrix object
attached to the LMVM Tao implementation. The history size must be
registered on the global PETSc options database before the
\texttt{TAOSolver} is constructed,
\begin{lstlisting}[language=Python]
PETSc.Options().setValue('-tao_lmvm_mat_lmvm_hist_size', 100)
\end{lstlisting}
or the entire pyadjoint TAO solver must be constructed with an
empty \texttt{options\_prefix}. We use the former in
\cref{sec:meshdep:implementation}.

\paragraph{Does this generalise to a non-quadratic problem?} The
Poisson optimal control problem is quadratic in the control $m$, and
quasi-Newton methods are known to be unusually effective on
quadratic problems regardless of the inner product used. To check
that the mesh-dependence pattern was not an artefact of the
quadratic structure, we repeated the experiment on a one-dimensional
viscous Burgers control problem, exposed through
\texttt{meshdep.problems.burgers\_control.solve\_burgers\_control}.
The forward equation is the backward-Euler discretisation of
\begin{equation}
  u_{t} + u\, u_{x} - \nu\, u_{xx} = m,
  \qquad x \in (0, 1),\ t \in (0, T),
\end{equation}
with $\nu = 10^{-3}$ and homogeneous Dirichlet boundary conditions.
The objective penalises misfit at the terminal time. We sweep the
same stretching parameter as on Poisson and run TAO/LMVM with the
default history of five on three configurations:
$L^{2}$ Riesz map with $\varepsilon = 10^{-7}$ ($L^{2}$);
$L^{2}$ Riesz map with $\varepsilon = 10^{-5}$ (loose $L^{2}$);
and $H^{1}$ Riesz map with $\varepsilon = 10^{-3}$ (fair $H^{1}$).

\begin{table}[H]
  \centering
  \begin{tabular}{l r r r}
    \toprule
    configuration
       & \multicolumn{3}{c}{realised $h_{\max} / h_{\min}$} \\
    \cmidrule(lr){2-4}
       & $1.0$ & $6.9$ & $18.3$ \\
    \midrule
    $L^{2}$,\, $\varepsilon=10^{-7}$, hist $5$ & $130$ & $116$ & $115$ \\
    $L^{2}$,\, $\varepsilon=10^{-5}$, hist $5$ & $51$ & $44$ & $41$ \\
    $H^{1}$,\, $\varepsilon=10^{-3}$, hist $5$ & $26$ & $30$ & $23$ \\
    \bottomrule
  \end{tabular}
  \caption{TAO/LMVM iteration counts on the one-dimensional viscous
  Burgers control problem with default history size five. None of
  the three configurations exhibits the mesh-dependent blow-up seen
  on the Poisson problem; instead the iteration count is essentially
  flat across the grading sweep.}
  \label{tab:meshdep:burgers}
\end{table}

The headline observation is that the dramatic blow-up of LMVM
iterations with grading that we documented on Poisson does
\emph{not} transfer to Burgers: the iteration count is essentially
flat across the grading sweep in all three configurations. This is
consistent with the diagnosis above. The graded-mesh Poisson
Hessian has a long-tailed spectrum because the discrete Laplacian
is very ill-conditioned on highly anisotropic grids; the
discrete Hessian of the time-dependent Burgers control problem is
better conditioned because the implicit time stepping acts as a
mass-matrix preconditioner on the high-frequency modes. The
mesh-dependence pattern of \citet[Table~2.2(a)]{shfp2017} is
therefore really a statement about the conditioning of the
discrete Hessian on graded meshes for stationary diffusion-dominated
problems, rather than a universal pathology of the $\ell^{2}$ inner
product.

\subsection{Discussion}
\label{sec:meshdep:discussion}

The Panel (a) results reproduced the central qualitative finding of
\citet[Sect.~2.3]{shfp2017}: when the natural function-space inner
product of the control is respected by the optimisation algorithm,
the iteration count of a quasi-Newton method is essentially
independent of the mesh ratio, while replacing it with the ambient
$\ell^{2}$ inner product causes the iteration count to grow at
roughly the polynomial rate predicted by the analysis in
\citet[Sect.~2.2.4]{shfp2017}.

Three observations from the implementation work bear on the
practical use of the optimisation packages catalogued by
\citet{shfp2017}. First, attaching a \texttt{riesz\_map} to a
pyadjoint \texttt{Control} \emph{is} enough for PETSc TAO/LMVM to
behave as a function-space quasi-Newton method in the algorithmic
sense - contrary to a conjecture that an earlier version of this
chapter recorded, the LMVM Mat does respect the inner product
induced by the initial Hessian preconditioner in its rank-two
updates. The original failure to reproduce the Moola row was
caused by the limited-memory approximation being insufficient,
not by an internal use of the $\ell^{2}$ inner product. The
diagnostic chain that led to this corrected picture is recorded in
\cref{sec:meshdep:tao-revisited}: with a sweep of the LMVM history
size at a tight tolerance, the final gradient norm on a graded mesh
falls from $10^{-11}$ at history five to machine precision at
history twenty-five, and on the Burgers control problem the
default history is already sufficient. So the practical recipe for
reproducing the Moola row via stock TAO/LMVM is to construct
\texttt{Control(m, riesz\_map="L2")} and to set
\texttt{tao\_lmvm\_mat\_lmvm\_hist\_size} to a value that captures
the relevant part of the discrete Hessian spectrum.

Second, although TAO/LMVM is in principle sufficient, the
hand-written Hilbert-space L-BFGS of
\cref{sec:meshdep:hilbert-lbfgs} remains useful as a reference
implementation: every dot product used in the two-loop recursion is
explicit in the Python source, making it straightforward to verify
where the chosen inner product enters the algorithm. The
hand-written implementation also avoids a subtlety of pyadjoint's
TAO wrapper that makes the LMVM history hard to set through the
public \texttt{parameters} interface, as documented in
\cref{sec:meshdep:tao-revisited}.

Third, the discrepancy between our Panel (b) results and those
reported by \citet[Table~2.2(b)]{shfp2017} indicates that the
discussion in the reference text does not fully specify how the
Moola row is computed in the $H^{1}$ case. The Hilbert-space
L-BFGS implementation converges to the correct minimum but takes
two orders of magnitude more iterations than the value reported in
the book; we attribute this primarily to the convergence threshold
being expressed in $H^{1}$, which is a stricter norm than $L^{2}$,
and partly to the high condition number of the discrete Hessian in
the $H^{1}$ inner product on the unit square. The same
\texttt{Control(m, riesz\_map="H1")} routed through TAO/LMVM with a
mesh-fair $H^{1}$ tolerance converges in roughly twenty-five
iterations across the grading sweep, indicating that the $H^{1}$
inner product itself does deliver mesh-independent behaviour;
identifying the four-iteration Moola step would require inspection
of the Moola source code, which is not currently possible in a
Firedrake environment.


% =============================================================================
\section{Conclusions and Future Work}
\label{sec:conclusions}
% =============================================================================

This work has developed the mathematical and numerical infrastructure
for adjoint-based optimal control of two model problems. We derived
the weak formulation, the adjoint equation and the $L^2$-gradient of
the reduced functional for the linear Poisson problem in
\cref{sec:poisson-weak,sec:poisson-optimisation,sec:poisson-derivative},
and verified the resulting numerical pipeline through manufactured
solutions and Taylor tests in \cref{sec:poisson-numerics}. We then
extended the same machinery to the nonlinear, time-dependent scalar
viscous Burgers problem in
\cref{sec:burgers-weak,sec:burgers-derivative,sec:burgers-numerics},
showing that the structure of the adjoint analysis carries over to a
problem whose adjoint must be solved backward in time.

\Cref{sec:mesh-dependence} reproduced Panel~(a) of Table~2.2 of
\citet{shfp2017} on six Netgen graded meshes spanning realised
ratios from $5.92$ to $193.88$. A custom implementation of
BFGS-in-Hilbert-space following \citet[Sect.~1.5.4.1]{shfp2017},
in which every dot product in the two-loop recursion is taken in
the $L^{2}$ inner product, produced an iteration count that varied
only between seven and nine across the entire range of mesh ratios.
The SciPy L-BFGS-B row, run on the ambient $\ell^{2}$ inner product
with the convergence criterion checked externally in $L^{2}$,
produced an iteration count that doubled at every refinement, in
agreement with the polynomial bound derived in
\citet[Sect.~2.2.4]{shfp2017}. Routing the same problem through
PETSc TAO/LMVM with \texttt{Control(m, riesz\_map="L2")} also
reproduces the qualitative pattern, provided that the LMVM history
size is enlarged from the PETSc default of five (which is too
small to capture the spectrum of the discrete Hessian on the
strongly graded meshes); the diagnostic chain that established
this, including a noise-floor probe at machine precision and a
Burgers control problem that does not exhibit the same
mesh-dependence, is documented in
\cref{sec:meshdep:tao-revisited}. Panel~(b) of Table~2.2 was only
partially reproduced; the open question of how Moola achieves a
four-iteration count in $H^{1}$ is discussed in
\cref{sec:meshdep:panel-b}.

Several directions remain open. The vector-valued time-dependent
Burgers problem treated by \citet{shfp2017} would be a natural
extension of \cref{sec:burgers-derivative}, and would let us
reproduce the mesh-dependence experiments of their Chapter~3 on
randomly refined meshes. A second direction is to use the
inner-product-aware optimisation framework on a problem with
practical relevance, for example tidal-turbine array layout
optimisation, as in \citet[Chapter~3]{shfp2017}. A third direction,
on the implementation side, is to extend the \texttt{meshdep} package
with Newton-CG and IPOPT wrappers so that the full grid of methods
considered by \citet{shfp2017} is reproducible.

A fourth direction, prompted by the partial reproduction of
\cref{sec:meshdep:panel-b}, is to identify the additional
algorithmic step that lets \citet{shfp2017}'s Moola row converge in
four iterations in the $H^{1}$ case; this would likely require
inspection of the Moola source code and the development of a
Firedrake-compatible build of that library, which is currently
blocked by Moola's hard-coded \texttt{dolfin} dependency.


% =============================================================================
\bibliographystyle{plainnat}
\bibliography{biblio}
% =============================================================================

\end{document}
